\DeclareMathOperator{\supp}{supp}
\DeclareMathOperator{\skewness}{Skewness}
\DeclareMathOperator{\DU}{D-U}
\DeclareMathOperator{\CU}{C-U}
\DeclareMathOperator{\Geo}{Geo}
\DeclareMathOperator{\Pois}{Pois}
\DeclareMathOperator{\Gam}{Gam}
\DeclareMathOperator{\LN}{LogN}
\DeclareMathOperator{\cov}{Cov}
\DeclareMathOperator*{\argmax}{arg\,max}
\DeclareMathOperator*{\argmin}{arg\,min}
\DeclareMathOperator{\re}{Re}
\DeclareMathOperator{\im}{Im}

\newcommand{\N}{\mathbb{N}}
\newcommand{\Z}{\mathbb{Z}}
\newcommand{\Q}{\mathbb{Q}}
\newcommand{\R}{\mathbb{R}}
\newcommand{\C}{\mathbb{C}}
\newcommand{\G}{\mathcal{G}}
\newcommand{\mP}{\mathbb{P}}
\newcommand{\F}{\mathcal{F}}
\newcommand{\B}{\mathcal{B}}
\newcommand{\1}{\mathbbm{1}}
\newcommand{\Poi}{\text{Poi}}
\newcommand{\Var}{\text{Var}}
\newcommand{\as}{\text{a.s.}}
\newcommand{\vect}[1]{\boldsymbol{\mathbf{#1}}}
\newcommand{\paren}[1]{\left({#1}\right)}
\newcommand{\brk}[1]{\left[{#1}\right]}
\newcommand{\curl}[1]{\left\{{#1}\right\}}
\providecommand{\lambdaup}{\uplambda}
\renewcommand{\lambda}{\lambdaup}
\providecommand{\wideparen}[1]{\overset{\frown}{#1}}

\theoremstyle{definition} % 定理環境 style 改用正體
\newtheorem{theorem}{定理}
\newtheorem*{theorem*}{定理}
\newtheorem{definition}{定義}
\newtheorem*{definition*}{定義}
\newtheorem{example}{例子}
\newtheorem*{example*}{例子}
\newtheorem{exercise}{練習}
\newtheorem*{exercise*}{練習}
\newtheorem{corollary}{推論}
\newtheorem*{corollary*}{推論}
\newtheorem{property}{性質}
\newtheorem*{property*}{性質}
\newtheorem{prop}{命題}
\newtheorem*{prop*}{命題}
\newtheorem{proposition}[prop]{命題}
\newtheorem*{proposition*}{命題}
\newtheorem{lemma}{引理}
\newtheorem*{lemma*}{引理}
\newtheorem{problem}{問題}
\newtheorem*{problem*}{問題}
\newtheorem*{remark}{註}
\newtheorem{discussion}{論例}
\newtheorem*{discussion*}{論例}
\newtheorem{homework}{作業}
\newtheorem*{homework*}{作業}
\newtheorem{example-non}{例}

% 讓 \begin{solution} 與 \begin{solution*} 皆可用，並以 # 或 \square 結尾
\newtheorem*{solution}{【解】}
\AtEndEnvironment{solution}{\hfill$\#$}

% 每章（\section）自動重置定理、例題、論例、命題、公式等計數器，確保全冊與單章編譯編號完全一致
\makeatletter
\@addtoreset{theorem}{section}
\@addtoreset{definition}{section}
\@addtoreset{example}{section}
\@addtoreset{exercise}{section}
\@addtoreset{corollary}{section}
\@addtoreset{property}{section}
\@addtoreset{prop}{section}
\@addtoreset{lemma}{section}
\@addtoreset{problem}{section}
\@addtoreset{discussion}{section}
\@addtoreset{homework}{section}
\@addtoreset{example-non}{section}
\@addtoreset{equation}{section}
\makeatother

% 當原書該章各節（§ X-Y）獨立重新起算定理、例題、論例、公式編號時，於 \subsection 後呼叫 \resetsectcounters
\newcommand{\resetsectcounters}{%
  \setcounter{theorem}{0}%
  \setcounter{definition}{0}%
  \setcounter{example}{0}%
  \setcounter{exercise}{0}%
  \setcounter{corollary}{0}%
  \setcounter{property}{0}%
  \setcounter{prop}{0}%
  \setcounter{lemma}{0}%
  \setcounter{problem}{0}%
  \setcounter{discussion}{0}%
  \setcounter{homework}{0}%
  \setcounter{example-non}{0}%
  \setcounter{equation}{0}%
}

\renewcommand{\figurename}{圖}
\renewcommand{\tablename}{表}

% 統一勘誤表返回連結巨集：各單章主檔與全冊主檔皆於勘誤表標題處定義 \label{sec:current_errata}
\providecommand{\errataref}{\hyperref[sec:current_errata]{目錄前勘誤表}}
